% §17.8 Bibliographic Notes
%
% Cross-reference resolution status:
%
% DEFINED here:
%   \label{sec:bib-notes-ch17}                -- this section
%
% INTERNAL refs (chapter 17):
%   \Cref{sec:what-is-cipher-map}             -- §17.1
%   \Cref{sec:fhe-sidebar}                    -- §17.4
%   \Cref{sec:boolean-trapdoor}               -- §17.5
%   \Cref{sec:maximizing-confidentiality}     -- §17.6
%
% Citations introduced here:
%   \cite{towell-cipher-maps}                 -- new bib entry
%   \cite{towell-boolean-trapdoor}            -- new bib entry
%   \cite{towell-maximizing-confidentiality}  -- new bib entry
%   \cite{gentry-2009}                        -- new bib entry
%   \cite{naor-yung-1989}                     -- new bib entry
%   \cite{goldreich-1996}                     -- already in bib from ch16
%
% Four paragraphs: cipher-maps manuscript; trapdoor-computing
% series (Boolean algebra over trapdoor sets, maximizing
% confidentiality, plus the algebraic cipher types thread);
% FHE literature (Gentry 2009); oblivious-computing tradition
% (Goldreich 1996, Naor-Yung 1989).

\section*{Bibliographic Notes}
\addcontentsline{toc}{section}{Bibliographic Notes}
\label{sec:bib-notes-ch17}

\paragraph{The cipher-maps manuscript.}
The cipher-map abstraction of \Cref{sec:what-is-cipher-map} is
developed at length in a companion manuscript,
\textcite{towell-cipher-maps}, ``Cipher maps: a unified framework
for oblivious function evaluation,'' which the present chapter
summarizes. The manuscript treats the abstraction in its own right
rather than as a special case of the Bernoulli framework, gives a
classification of cipher-map families by their channel-matrix
structure, and develops the cryptographic-side analysis (key
schedules, side-channel considerations, computational hardness
assumptions) that the present chapter has not entered. The
chapter's contribution beyond the manuscript is the identification
of cipher-map error rates with Bernoulli atoms and the lifting of
chapter~4's Kronecker composition to cipher-map composition; the
manuscript treats composition algebraically without anchoring it to
the Bernoulli square.

\paragraph{The trapdoor-computing series.}
The trapdoor-computing paradigm is the subject of an ongoing series
of which three drafts are directly relevant. The Boolean-algebra
section (\Cref{sec:boolean-trapdoor}) draws on
\textcite{towell-boolean-trapdoor}, ``Boolean algebra over trapdoor
sets,'' which gives the full algebraic treatment of union,
intersection, complement, and conditional operations over trapdoor
representations, including the cancellation phenomena the
present chapter alluded to. The maximizing-confidentiality section
(\Cref{sec:maximizing-confidentiality}) draws on
\textcite{towell-maximizing-confidentiality}, ``Maximizing
confidentiality in encrypted search,'' which develops the
PPV-as-confidentiality connection in the encrypted-search setting
and extends the analysis to information-theoretic measures
(observed-versus-maximum entropy ratios, compression-based entropy
estimation) the present chapter has not discussed. A third strand,
on algebraic cipher types, generalizes the chapter's set-and-map
treatment to algebraic data types and remains in active drafting.

\paragraph{Fully homomorphic encryption.}
The FHE sidebar of \Cref{sec:fhe-sidebar} stands on
\textcite{gentry-2009}, ``Fully Homomorphic Encryption Using Ideal
Lattices,'' which gave the first plausibly-secure fully
homomorphic scheme by combining a somewhat-homomorphic scheme with
a bootstrapping construction that refreshes ciphertexts faster
than the noise from homomorphic operations accumulates. The
scheme is conceptually clean and operationally heavy; modern
schemes (BFV, BGV, CKKS) reduce the per-operation cost by orders
of magnitude but inherit Gentry's bootstrapping idea and remain
several orders of magnitude slower than the cipher-map approach
that the present chapter develops. The cryptographic standard
references on FHE are surveys rather than primary sources, and
the reader interested in modern schemes should follow the
literature trail from Gentry's paper rather than from the
chapter's prose.

\paragraph{The oblivious-computing tradition.}
The trapdoor-paradigm framing of \Cref{sec:fhe-sidebar} sits in a
literature that goes back to early work on secure computation.
\textcite{goldreich-1996}, ``Foundations of Cryptography:
Fragments of a Book'' (already cited in chapter~16 for the
oblivious-Turing-machine model), gives the foundational treatment
of computation under cryptographic constraints, including the
relation between oblivious computation and the secure-multiparty
setting that cipher maps generalize. \textcite{naor-yung-1989},
``Universal One-Way Hash Functions and their Cryptographic
Applications,'' is the older entry point to the trapdoor-function
line of work that the cipher-map abstraction inherits its name
from. The chapter has used ``trapdoor'' in the colloquial sense
rather than the strict cryptographic sense, but the connection to
the formal trapdoor-function literature is real and the
encrypted-search applications of chapter~18 will activate it more
explicitly. The bibliographic-notes treatment is necessarily
brief; the secure-multiparty-computation literature is enormous,
and the reader seeking the full context should consult one of
the standard textbook references in cryptography.
